\section*{Bab \ref{ch:b2:metric} --- Topologi Ruang Metrik}

\begin{solution}{exo:b2:metric:1}
$\delta$: kesetangkupan dan pemisahannya jelas; untuk ketaksamaan
segitiganya: $\min(1, u + v) \leq \min(1,u) + \min(1,v)$ bagi $u, v \geq
0$ (jika salah satu minimumnya $1$, ruas kanannya $\geq 1$; jika tidak,
ruas kanannya $u + v$). Untuk $d$: ia tarikan balik $\abs{\cdot}$ oleh
$\arctan$ yang injektif, jadi ketiga aksiomanya berpindah.

Kekonvergenannya: untuk $\delta$, $\delta(x_n, x) \to 0 \iff \abs{x_n -
x} \to 0$ (sebab untuk nilai kecil kedua jaraknya berimpit), jadi barisan
konvergennya sama seperti biasa. Untuk $d$: $d(x_n, x) \to 0 \iff \arctan
x_n \to \arctan x \iff x_n \to x$ (lewat kekontinuan dan kemonotonan
tegas $\arctan$ beserta inversnya pada jangkauan yang bersangkutan), jadi
lagi-lagi kekonvergenan yang biasa.

Ruang $(\R, d)$ \emph{tidak} \omterm{def:b2:metric:complete}{lengkap}: $x_n = n$ memenuhi $d(x_p, x_q)
= \abs{\arctan p - \arctan q} \to 0$ (sebab keduanya menuju $\frac\pi2$),
jadi Cauchy; tetapi $(x_n)$ tidak konvergen untuk $d$ (sebab limitnya
terhadap $d$ akan menjadi limit yang biasa juga). Kelengkapan adalah
sifat \emph{jaraknya}, bukan sekadar sifat barisan konvergennya.
\end{solution}

\begin{solution}{exo:b2:metric:2}
Konvergen $\Rightarrow$ Cauchy: sebab $d(x_p, x_q) \leq d(x_p, \ell) +
d(\ell, x_q)$. Terbatas: setelah $N$, berlaku $d(x_n, \ell) \leq 1$;
sedangkan berhingga suku pertamanya pun berada di dalam suatu jari-jari.

Cauchy $+$ barisan bagian konvergen $x_{\varphi(n)} \to \ell$: diberikan
$\varepsilon$, untuk $n$ yang besar berlaku $d(x_n, \ell) \leq d(x_n,
x_{\varphi(n)}) + d(x_{\varphi(n)}, \ell) \leq 2\varepsilon$ (suku
pertamanya berkat sifat Cauchy, sebab $\varphi(n) \geq n$).

\omterm{def:b2:metric:compact}{Kompak} $\Rightarrow$ \omterm{def:b2:metric:complete}{lengkap}: barisan Cauchy punya barisan bagian yang
konvergen (kekompakan), jadi ia konvergen.
\end{solution}

\begin{solution}{exo:b2:metric:3}
$d_\infty(f, g) = \sup_{\intcc{0}{1}} \abs{x - x^2} = \frac14$
(sebab maksimum $x - x^2$ tercapai di $x = \frac12$).

$\overline B(0, 1) = \{f : \sup\abs f \leq 1\}$: yakni semua fungsi
\omterm{def:b2:metric:continuity}{kontinu} yang bernilai di $\intcc{-1}{1}$.

Himpunan $\{f : f(0) = 0\}$ adalah prapeta $\{0\}$ di bawah
\emph{evaluasi} $f \mapsto f(0)$, yang bersifat \omterm{def:b2:metric:continuity}{Lipschitz} berkonstanta
$1$ (sebab $\abs{f(0) - g(0)} \leq d_\infty(f,g)$), jadi \omterm{def:b2:metric:continuity}{kontinu}:
karenanya himpunannya tertutup
(\cref{thm:b2:metric:globalcontinuity}). Demikian pula $\{f :
f(0) > 0\}$ adalah prapeta $\intoo{0}{+\infty}$ yang \omterm{def:b2:metric:topology}{terbuka}: jadi \omterm{def:b2:metric:topology}{terbuka}.
\end{solution}

\begin{solution}{exo:b2:metric:4}
\emph{\omterm{def:b2:metric:complete}{Lengkap}:} barisan Cauchy dengan $\varepsilon = \frac12$ akhirnya
menjadi konstan, jadi ia konvergen.

\emph{\omterm{def:b2:metric:compact}{Kompak} bila dan hanya bila hingga:} jika $X$ hingga, sembarang
barisan mengambil suatu nilai tak hingga kali (jadi ada barisan bagian
yang konstan). Jika $X$ tak hingga, barisan berisi titik yang berbeda
sepasang demi sepasang punya semua jarak antara sama dengan $1$: jadi tak
ada barisan bagian yang Cauchy, sehingga tak ada yang konvergen.

\emph{Himpunan bagian yang \omterm{def:b2:metric:connected}{terhubung}:} yaitu himpunan tunggal (dan
$\emptyset$). Sembarang $A$ yang punya dua titik $x \neq y$ terpecah
sebagai $\{x\} \cup (A \setminus\{x\})$, dan keduanya \omterm{def:b2:metric:topology}{terbuka} di $A$
(sebab setiap himpunan bagian ruang diskret bersifat \omterm{def:b2:metric:topology}{terbuka} --- bola
berjari-jari $\frac12$ hanyalah himpunan tunggal).
\end{solution}

\begin{solution}{exo:b2:metric:5}
Fungsi $g(x) = d(x, f(x))$ bersifat \omterm{def:b2:metric:continuity}{kontinu} pada $X$ yang \omterm{def:b2:metric:compact}{kompak}
(sebab $\abs{g(x) - g(y)} \leq 2d(x,y)$ lewat dua ketaksamaan segitiga),
jadi ia mencapai minimumnya di suatu $a$
(\cref{thm:b2:metric:compactprops}). Jika $f(a) \neq a$, maka
\[
g\bigl(f(a)\bigr) = d\bigl(f(a), f(f(a))\bigr) < d(a, f(a)) = g(a),
\]
dan itu bertentangan dengan keminimalannya. Jadi $f(a) = a$;
ketunggalannya seperti biasa (dua titik tetap $a \neq b$ memberi $d(a,b)
= d(f(a), f(b)) < d(a,b)$).

Contoh tanpa kekompakan: $f(x) = x + \frac1x$ pada $\intco{1}{+\infty}$:
di sini $\abs{f(x) - f(y)} = \abs{x - y}\,\abs{1 - \frac{1}{xy}} < \abs{x
- y}$ untuk $x \neq y$ (sebab $xy > 1$), namun $f(x) > x$ di mana-mana.
\end{solution}

\begin{solution}{exo:b2:metric:6}
Pilih $x_n \in K_n$. Semua sukunya mulai peringkat $n$ berada di $K_n$;
khususnya seluruh barisannya berada di $K_0$ yang \omterm{def:b2:metric:compact}{kompak}: jadi sarikan
$x_{\varphi(k)} \to \ell$. Untuk tiap $n$ yang tetap, suku
$x_{\varphi(k)}$ dengan $\varphi(k) \geq n$ berada di $K_n$ yang
\emph{tertutup}, jadi limitnya $\ell \in K_n$. Karenanya $\ell \in
\bigcap K_n$.

Contoh penyangkal yang tertutup: $F_n = \intco{n}{+\infty}$ di $\R$:
bersarang, tertutup, tak kosong, tetapi irisannya kosong.
\end{solution}

\begin{solution}{exo:b2:metric:7}
Kekontinuan $f^{-1}$ berarti: peta $f(F)$ atas himpunan tertutup $F
\subseteq K$ bersifat tertutup (sebab prapeta di bawah $f^{-1}$ adalah
peta di bawah $f$). Himpunan tertutup $F$ di dalam $K$ yang \omterm{def:b2:metric:compact}{kompak}
bersifat \omterm{def:b2:metric:compact}{kompak} (\cref{thm:b2:metric:compactprops}~(1)); peta \omterm{def:b2:metric:continuity}{kontinunya}
$f(F)$ bersifat \omterm{def:b2:metric:compact}{kompak}, jadi tertutup. Karenanya $f^{-1}$ \omterm{def:b2:metric:continuity}{kontinu},
sehingga $f$ merupakan homeomorfisma.

Contoh penyangkal: $\gamma(t) = (\cos t, \sin t)$ dari
$\intco{0}{2\pi}$ (yang tidak \omterm{def:b2:metric:compact}{kompak}) pada lingkaran satuan merupakan
bijeksi \omterm{def:b2:metric:continuity}{kontinu}, tetapi $\gamma^{-1}$ tidak \omterm{def:b2:metric:continuity}{kontinu} di $(1,0)$:
sebab titik lingkaran yang tepat di bawah sumbunya berparameter dekat
$2\pi$, bukan dekat $0$.
\end{solution}

\begin{solution}{exo:b2:metric:8}
Di sini $C = \bigcap_n C_n$ dengan tiap $C_n$ (gabungan $2^n$ interval
tertutup berpanjang $3^{-n}$) bersifat tertutup: jadi $C$ tertutup dan
terbatas di $\R$, sehingga \omterm{def:b2:metric:compact}{kompak} (\cref{thm:b2:metric:compactprops}~(2)).

Interior kosong: $C$ tidak memuat interval berpanjang $> 3^{-n}$ (sebab
ia berada di dalam $C_n$, yang komponennya berpanjang segitu), untuk
setiap $n$.

Kardinalitasnya: titik $C$ persis semua bilangan real $\sum_{n\geq1}
a_n 3^{-n}$ berangka $a_n \in \{0, 2\}$ (sebab pada tiap tahapnya,
sepertiga tengah yang dibuang menyingkirkan angka $1$); dan pemetaan
$(a_n) \mapsto \sum a_n 3^{-n}$ merupakan bijeksi dari $\{0,2\}^{\N^*}$
pada $C$ (keinjektifannya seperti pada \cref{exo:b2:structures:3}). Jadi
$C$ \omterm{def:b2:structures:countable}{ekuipoten} dengan $\{0,1\}^{\N}$: tak \omterm{def:b2:structures:countable}{terbilang}, walaupun ``berpanjang nol''.
\end{solution}

\begin{solution}{exo:b2:metric:9}
Andaikan $a \notin f(X)$. Peta $f(X)$ bersifat \omterm{def:b2:metric:compact}{kompak} (sebab peta
\omterm{def:b2:metric:continuity}{kontinu}), jadi tertutup; karenanya
\[
\varepsilon = d\bigl(a, f(X)\bigr)
= \inf_{y \in f(X)} d(a, y) > 0
\]
(sebab infimum fungsi \omterm{def:b2:metric:continuity}{kontinu} pada himpunan \omterm{def:b2:metric:compact}{kompak} tercapai; dan
seandainya nilainya $0$, maka $a$ akan melekat pada $f(X)$ yang tertutup,
jadi berada di dalamnya).

Tinjaulah orbit $x_n = f^n(a)$ (dengan $x_0 = a$). Untuk $p < q$:
\[
d(x_p, x_q) = d\bigl(f^p(a), f^p(f^{q-p}(a))\bigr)
= d\bigl(a, f^{q-p}(a)\bigr) \geq \varepsilon,
\]
(sebab isometrinya diiterasikan $p$ kali; dan $f^{q-p}(a) \in f(X)$
karena $q - p \geq 1$). Barisan yang jarak antaranya $\geq \varepsilon$
tidak punya barisan bagian yang konvergen --- dan itu bertentangan
dengan kekompakan. Karenanya $f(X) = X$.
\end{solution}

\begin{solution}{exo:b2:metric:10}
Misalkan $A, B \in GL_n(\C)$ dan $p(z) = \det\bigl((1-z)A + zB\bigr)$:
sebuah polinomial dalam $z$ (sebab tiap unsurnya afin dalam $z$, dan
\omterm{def:b2:linalg:det}{determinannya} polinomial dalam unsurnya). Karena $p(0) = \det A \neq 0$,
$p$ tidak nol secara identik, jadi ia punya berhingga banyak akar $z_1,
\dots, z_m$ (dan tak satu pun sama dengan $0$ atau $1$, sebab $p(1) =
\det B \neq 0$). Bidang $\C$ dikurangi berhingga banyak titik bersifat
\omterm{def:b2:metric:connected}{terhubung} lintasan: jadi ada lintasan dari $0$ ke $1$ yang menghindari
$z_i$ (ambil garis patah lewat sebuah titik yang jauh dari semua akarnya,
atau sebuah busur lingkaran; halangannya hanya berhingga banyak). Di
sepanjang lintasan $\gamma$ semacam itu, $t \mapsto (1 - \gamma(t))A +
\gamma(t)B$ merupakan lintasan \omterm{def:b2:metric:continuity}{kontinu} \emph{di dalam} $GL_n(\C)$ dari
$A$ ke $B$ (sebab \omterm{def:b2:linalg:det}{determinannya} tak pernah nol di sana). Karenanya
$GL_n(\C)$ \omterm{def:b2:metric:connected}{terhubung} lintasan --- tidak seperti sepupu realnya
(\cref{ex:b2:metric:glnr}): bidang kompleks punya ruang untuk berjalan
memutari halangan.
\end{solution}

\begin{solution}{exo:b2:metric:11}
\emph{\omterm{def:b2:metric:continuity}{Lipschitz}:} untuk $a \in A$ berlaku $d(x, a) \leq d(x, y) + d(y,
a)$; ambil infimumnya atas $a$: $d(x, A) \leq d(x, y) + d(y,
A)$, lalu tukarkan $x, y$: $\abs{d(x,A) - d(y,A)} \leq d(x,y)$.

\emph{Penolkannya:} $d(x, A) = 0$ bila dan hanya bila ada $a_n \in A$
dengan $d(x, a_n) \to 0$, bila dan hanya bila $x$ limit titik $A$, bila
dan hanya bila $x \in \overline A$.

\emph{Sakelarnya:} untuk $A, B$ tertutup yang saling lepas: penyebut
$d(x,A) + d(x,B)$ tak pernah nol (sebab itu akan memaksa $x \in \overline
A \cap \overline B = A \cap B = \emptyset$), jadi $\varphi$ terdefinisi
dengan baik dan \omterm{def:b2:metric:continuity}{kontinu} sebagai hasil bagi fungsi \omterm{def:b2:metric:continuity}{kontinu} berpenyebut tak
nol. Lalu $\varphi(x) = 0$ bila dan hanya bila $d(x, A) = 0$, bila dan
hanya bila $x \in A$; dan $\varphi(x) = 1$ bila dan hanya bila $d(x, B) =
0$, bila dan hanya bila $x \in B$; sedangkan $0 \leq \varphi \leq 1$ di mana-mana.
\end{solution}

\begin{solution}{exo:b2:metric:12}
Misalkan $B(x_0, r_0)$ sembarang bola; kita cari titik $\bigcap U_n$
di dalamnya. Karena $U_1$ padat dan \omterm{def:b2:metric:topology}{terbuka}, $U_1 \cap B(x_0, r_0)$
tak kosong dan \omterm{def:b2:metric:topology}{terbuka}: jadi ia memuat bola tertutup $\overline B(x_1,
r_1)$ dengan $0 < r_1 \leq \frac{r_0}2$ (perkecil saja jari-jarinya).
Secara induktif, $U_{n+1} \cap B(x_n, r_n)$ tak kosong dan \omterm{def:b2:metric:topology}{terbuka}: pilih
$\overline B(x_{n+1}, r_{n+1}) \subseteq U_{n+1} \cap B(x_n,
r_n)$ dengan $r_{n+1} \leq \frac{r_n}2$. Untuk $p, q \geq n$, kedua titik
$x_p, x_q$ berada di $B(x_n, r_n)$ dengan $r_n \leq 2^{-n}r_0$: jadi
barisannya Cauchy, sehingga ia konvergen ke suatu $\ell$ berkat
kelengkapannya. Untuk tiap $n$, ekor barisannya berada di bola
\emph{tertutup} $\overline B(x_{n+1}, r_{n+1}) \subseteq U_{n+1}
\cap B(x_0, r_0)$, jadi $\ell \in U_{n+1}$ untuk setiap $n$ dan $\ell
\in \overline B(x_1, r_1) \subseteq B(x_0, r_0)$. Karenanya
$\bigcap_n U_n$ bertemu setiap bola: padat.

\emph{Penerapannya:} jika $\R = \bigcup_n F_n$ dengan tiap $F_n$
tertutup dan berinterior kosong, maka $U_n = \R \setminus F_n$ bersifat
padat (sebab $\overline{U_n} = \R$ bila dan hanya bila $F_n$ berinterior
kosong) dan \omterm{def:b2:metric:topology}{terbuka}, sehingga Baire memberi sebuah titik di $\bigcap U_n
= \R \setminus \bigcup F_n$: kontradiksi. Khususnya $\R \neq \bigcup_{x
\in D} \{x\}$ untuk $D$ yang \omterm{def:b2:structures:countable}{terbilang} (sebab himpunan tunggal bersifat
tertutup dan berinterior kosong): jadi $\R$ tidak \omterm{def:b2:structures:countable}{terbilang} ---
\cref{thm:b2:structures:cantor} lewat jalan lain.
\end{solution}

\begin{solution}{pb:b2:metric:1}
\textbf{1.} Misalkan $F$ tertutup di $X$ yang \omterm{def:b2:metric:complete}{lengkap} dan $(x_n)
\subseteq F$ bersifat Cauchy: ia konvergen di $X$ ke suatu $\ell$, dan
$\ell \in F$ karena $F$ tertutup (sebab limit barisan di $F$
tetap di $\overline F = F$): jadi $F$ \omterm{def:b2:metric:complete}{lengkap}. Sebaliknya misalkan $A
\subseteq X$ \omterm{def:b2:metric:complete}{lengkap} dan $x \in \overline A$: ada barisan
di $A$ yang konvergen ke $x$; ia Cauchy, jadi ia konvergen \emph{di
dalam} $A$; karena limitnya tunggal, $x \in A$, sehingga $A$ tertutup.
Karena $\bigl(C(I), d_\infty\bigr)$ bersifat \omterm{def:b2:metric:complete}{lengkap}
(\cref{thm:b2:metric:rncomplete}), himpunan bagiannya yang tertutup
bersifat \omterm{def:b2:metric:complete}{lengkap}.

\textbf{2.} Untuk $y, z \in C(I)$ dan $t \in I$:
\[
\Bigl|\int_{t_0}^{t}y - \int_{t_0}^{t}z\Bigr| \leq
\abs{t - t_0}\,\sup_I\abs{y - z} \leq h\,d_\infty(y, z),
\]
lalu ambil supremumnya atas $t$.

\textbf{3.} Dengan memakai kedua persamaan titik tetapnya dan
ketaksamaan segitiga:
\[
d(\ell_g, \ell_{\widetilde g})
= d\bigl(g(\ell_g), \widetilde g(\ell_{\widetilde g})\bigr)
\leq d\bigl(g(\ell_g), g(\ell_{\widetilde g})\bigr) +
d\bigl(g(\ell_{\widetilde g}), \widetilde g(\ell_{\widetilde
g})\bigr)
\leq k\,d(\ell_g, \ell_{\widetilde g}) +
d\bigl(g(\ell_{\widetilde g}), \widetilde g(\ell_{\widetilde
g})\bigr),
\]
lalu selesaikan untuk $d(\ell_g, \ell_{\widetilde g})$ (sebab koefisien
$1 - k$ positif). Ketaksamaan keduanya membatasi jurang yang dinilai itu
oleh jurang yang seragam.

\textbf{4.} $g^m$ merupakan kontraksi pada ruang \omterm{def:b2:metric:complete}{lengkap} yang tak kosong,
jadi ia punya tepat satu titik tetap $\ell$
(\cref{thm:b2:metric:banach}). Maka $g^m(g(\ell)) =
g(g^m(\ell)) = g(\ell)$: jadi $g(\ell)$ titik tetap $g^m$,
sehingga $g(\ell) = \ell$ berkat ketunggalannya. Sebaliknya setiap titik
tetap $g$ juga titik tetap $g^m$: jadi ketunggalan bagi $g$. Orbitnya:
tetapkan $r \in \{0, \dots, m-1\}$; barisan bagian $(x_{qm + r})_q$
adalah orbit $g^m$ yang berawal di $x_r$, jadi ia konvergen ke $\ell$
ketika $q \to \infty$ (lewat Banach lagi). Kesepuluh
--- lebih tepatnya semua $m$ --- barisan bagiannya konvergen ke $\ell$
yang sama, jadi $x_n \to \ell$: diberikan $\varepsilon$, tiap kelas
sisanya akhirnya berada dalam $\varepsilon$, dan kelasnya berhingga banyak.

\textbf{5.} Jika $y$ \omterm{def:b2:metric:continuity}{kontinu} dan bernilai di $\intcc{y_0 -
b}{y_0 + b}$, maka integrannya $s \mapsto f(s, y(s))$ \omterm{def:b2:metric:continuity}{kontinu}
pada $I$ (sebab susunan fungsi \omterm{def:b2:metric:continuity}{kontinu}), sehingga ruas kanannya bersifat
$C^1$ dengan turunan $f(t, y(t))$ (teorema dasar kalkulus, jilid Tahun
ke-1). Jika $y$ memenuhi persamaan integralnya, maka ia fungsi
$C^1$ tersebut, $y(t_0) = y_0$, dan $y' = f(t, y)$. Sebaliknya,
mengintegralkan $y' = f(s, y(s))$ dari $t_0$ sampai $t$ memberi
persamaan integral itu.

\textbf{6.} Himpunan $X_h$ memuat fungsi konstan $y_0$; ia tertutup
sebagai prapeta $\intcc{0}{b}$ di bawah pemetaan \omterm{def:b2:metric:continuity}{kontinu} $y
\mapsto d_\infty(y, y_0)$ (sebab jarak bersifat \omterm{def:b2:metric:continuity}{Lipschitz} berkonstanta
$1$), jadi ia \omterm{def:b2:metric:complete}{lengkap} menurut pertanyaan 1. Kestabilannya: untuk $y \in
X_h$ dan $t \in I$,
\[
\abs{T(y)(t) - y_0} = \Bigl|\int_{t_0}^{t}f(s, y(s))\dd s\Bigr|
\leq M\abs{t - t_0} \leq Mh \leq b ,
\]
langkah terakhirnya berkat $h \leq b/M$ (atau $M = 0$, yang sepele). Dan
$T(y)$ bersifat \omterm{def:b2:metric:continuity}{kontinu} (bahkan $C^1$, pertanyaan 5): jadi $T(y) \in X_h$.

\textbf{7.} Untuk $t \in I$:
\[
\abs{T(y)(t) - T(z)(t)} \leq \int_{t_0}^{t}\abs{f(s, y(s)) -
f(s, z(s))}\,\abs{\dd s} \leq L\abs{t - t_0}\,d_\infty(y,z)
\leq Lh\,d_\infty(y,z).
\]
Jika $Lh < 1$: maka $T$ merupakan kontraksi pada $X_h$ yang \omterm{def:b2:metric:complete}{lengkap} dan
tak kosong, sehingga Banach memberi tepat satu titik tetap --- yang
menurut pertanyaan 5 adalah satu-satunya penyelesaiannya.

\textbf{8.} Secara induktif; kasus $n = 1$ tak lain ketaksamaan tengah
pada pertanyaan 7. Dengan mengandaikan batasnya untuk $n$, maka untuk $t
\geq t_0$ (kasus $t \leq t_0$ setangkup):
\[
\abs{T^{n+1}(y)(t) - T^{n+1}(z)(t)}
\leq L\int_{t_0}^{t}\abs{T^n(y)(s) - T^n(z)(s)}\dd s
\leq L\int_{t_0}^{t}\frac{L^n(s - t_0)^n}{n!}\dd s\;
d_\infty(y,z),
\]
dan integralnya bernilai $\frac{L^n(t -
t_0)^{n+1}}{(n+1)!}$: itulah batas untuk $n + 1$. Karenanya
$d_\infty(T^n y, T^n z) \leq \frac{(Lh)^n}{n!}d_\infty(y, z)$,
dan $\frac{(Lh)^n}{n!} \to 0$ (sebab deret eksponensialnya
konvergen): jadi suatu $T^m$ merupakan kontraksi. Pertanyaan 4 berlaku
pada $X_h$ yang \omterm{def:b2:metric:complete}{lengkap}: $T$ punya tepat satu titik tetap, yakni
masalah Cauchy itu punya tepat satu penyelesaian pada $I$ yang bernilai
di $\intcc{y_0 - b}{y_0 + b}$.

\textbf{9.} Misalkan $y$ menyelesaikan masalahnya pada $I$ dan andaikan
himpunan $E = \{t \in I, t > t_0 : \abs{y(t) - y_0} > b\}$ tak kosong
(sisi $t < t_0$ setangkup); misalkan $\tau = \inf E$. Berkat
kekontinuannya, $\abs{y(s) - y_0} \leq b$ untuk $s \in
\intcc{t_0}{\tau}$, jadi grafiknya berada di $R$ di sana, persamaan
integralnya berlaku sampai $\tau$, dan
\[
\abs{y(\tau) - y_0} = \Bigl|\int_{t_0}^{\tau}f\bigl(s,
y(s)\bigr)\dd s\Bigr| \leq M(\tau - t_0) \leq Mh \leq b .
\]
Jika $\tau < t_0 + h$, maka titik $E$ yang sedekat-dekatnya ke $\tau$
dari kanan memberi, berkat kekontinuan, $\abs{y(\tau) - y_0} \geq
b$, jadi $= b$; tetapi ungkapan di atas lalu memaksa $M(\tau - t_0) =
Mh$, yakni $\tau = t_0 + h$: kontradiksi. Jadi $\tau = t_0 + h$,
sehingga $E \subseteq \{t_0 + h\}$, dan ungkapan tadi (di $\tau = t_0 +
h$) memberi $\abs{y(t_0 + h) - y_0} \leq b$, yang bertentangan dengan
keanggotaannya di $E$. Karenanya $E = \emptyset$: setiap penyelesaian
pada $I$ tetap di dalam pitanya, merupakan titik tetap $T$ di $X_h$, dan
ketunggalannya berlaku tanpa syarat.

\textbf{10.} Di sini $T(y)(t) = 1 + \int_0^t y$. Dari $y^{(0)} \equiv
1$:
\[
y^{(1)}(t) = 1 + t,\quad
y^{(2)}(t) = 1 + t + \frac{t^2}2,\quad\dots\quad
y^{(n)}(t) = \sum_{k=0}^{n}\frac{t^k}{k!}
\]
(secara induktif: mengintegralkan jumlah parsialnya menambahkan suku
berikutnya). Inilah jumlah parsial Taylor bagi $\eu^{\,t}$; pada
sembarang $I$ yang terbatas keduanya konvergen seragam ke $\eu^{\,t}$
(sebab ekornya didominasi deret numerik konvergen $\sum h^k/k!$),
dan itu memang satu-satunya penyelesaiannya.

\textbf{11.} Fungsi $y \equiv 0$ merupakan penyelesaian. Untuk $y_c$: ia
bersifat $C^1$ (kedua potongannya demikian, dan di $t = c$ turunannya
cocok: $0$ dan $2(t - c) \to 0$), lalu untuk $t > c$: $y_c' = 2(t - c) =
2\sqrt{(t-c)^2} = 2\sqrt{\abs{y_c}}$; sedangkan untuk $t \leq c$ kedua
ruasnya nol. Jadi masalah Cauchy $y(0) = 0$ punya tak hingga banyak
penyelesaian ($c \geq 0$ sembarang, ditambah $y \equiv 0$). Medannya
$\varphi(y) = 2\sqrt{\abs y}$ tidak \omterm{def:b2:metric:continuity}{Lipschitz} di dekat $0$, sebab
$\frac{\varphi(y) - \varphi(0)}{y - 0} = \frac{2}{\sqrt y} \to
+\infty$ ketika $y \to 0^+$: jadi tak ada konstanta $L$ yang berhasil
pada persekitaran $0$ mana pun --- persis di tempat semua penyelesaiannya
bercabang.

\textbf{12.} Dengan memisahkan peubahnya (atau memeriksanya langsung),
satu-satunya penyelesaian lokalnya adalah $y(t) = \frac1{1 - t}$, yang
terdefinisi pada $\intoo{-\infty}{1}$ dan meledak di $t = 1$. Untuk
persegi panjang $\intcc{-a}{a} \times \intcc{1 - b}{1 + b}$: $M =
\sup y^2 = (1 + b)^2$, jadi setengah lebar yang tersahkan adalah $h =
\min\bigl(a, \frac{b}{(1+b)^2}\bigr)$. Memaksimalkan
$\frac{b}{(1+b)^2}$: turunannya nol di $b = 1$, dengan nilai
$\frac14$. Jadi teoremanya hanya menjamin kehidupan pada
$\intcc{-\frac14}{\frac14}$ --- yang memang lebih pendek daripada masa
hidup sejatinya, yaitu $1$ ke depan, dan jauh lebih pendek ke belakang:
teoremanya bersifat lokal menurut kodratnya, dan ledakan itu menunjukkan
ia tak mungkin lain.

\textbf{13.} Pemetaan $g$ mengirim $\intcc12$ ke dirinya sendiri: $g$
turun pada $\intcc{1}{\sqrt2}$ lalu naik sesudahnya (telaahlah $g'(x) =
\frac12 - \frac1{x^2}$), dengan $g(1) = g(2) = \frac32$ dan minimum
$g(\sqrt2) = \sqrt2 > 1$: jadi $g(\intcc12) \subseteq
\intcc{\sqrt2}{\frac32} \subseteq \intcc12$; dan $g$ mengirim
bilangan rasional ke bilangan rasional. Kontraksinya: $\abs{g'(x)} =
\abs{\frac12 - \frac1{x^2}} \leq \frac12$ pada $\intcc12$ (sebab
$\frac1{x^2} \in \intcc{\frac14}{1}$), jadi ketaksamaan nilai rata-rata
memberi $\abs{g(x) - g(y)} \leq \frac12\abs{x - y}$. Titik tetapnya
memenuhi $\frac x2 = \frac1x$, yakni $x^2 = 2$: mustahil di
$\Q$. Hipotesis yang gagal adalah kelengkapan $X$ (sebab $\Q \cap
\intcc12$ tidak \omterm{def:b2:metric:complete}{lengkap}); di dalam pelengkapannya, yaitu $\intcc12$,
titik tetapnya adalah $\sqrt2$ --- teorema Banach yang dijalankan pada
bilangan rasional justru \emph{menciptakan} bilangan irasionalnya.

\textbf{14.} Aposteriori: $d(x_n, \ell) \leq d(x_n, x_{n+1}) +
d(x_{n+1}, \ell) \leq k\,d(x_{n-1}, x_n) + k\,d(x_n, \ell)$,
sehingga $d(x_n, \ell) \leq \frac{k}{1-k}d(x_n, x_{n-1})$. Heron
dari $x_0 = \frac32$: $x_1 = \frac{17}{12}$, $d(x_1, x_0) =
\frac1{12}$, $k = \frac12$: jadi batas apriori
$\frac{k^n}{1-k}d(x_1,x_0) = \frac{2^{-n+1}}{12}$ baru turun di bawah
$10^{-6}$ pada $n = 18$. Kenyataannya $x_1 = \frac{17}{12}
\approx 1.41667$ (galat $2.5\cdot10^{-3}$), $x_2 =
\frac{577}{408} \approx 1.4142157$ (galat $2.1\cdot10^{-6}$),
$x_3 \approx 1.41421356237469$ (galat $1.6\cdot10^{-12}$):
jadi tiga langkah sudah cukup. Tiap langkah Heron kurang lebih
\emph{mengkuadratkan} galatnya (kekonvergenan kuadratik, gejala Newton:
\cref{ch:b2:realfun}); sedangkan taksiran kontraksinya, yang hanya
memaruhkannya, jujur untuk kasus terburuk tetapi pesimistis di sini.

\textbf{15.} Terapkan pertanyaan 3 dengan $g = T_{y_0}$ (sebuah
kontraksi-$Lh$, dengan $Lh < 1$) dan $\widetilde g = T_{z_0}$, yang
titik tetapnya $y[z_0]$. Untuk sembarang $y$,
\[
\abs{T_{y_0}(y)(t) - T_{z_0}(y)(t)} = \abs{y_0 - z_0},
\]
(sebab integralnya sama persis), jadi $\sup_y
d_\infty(T_{y_0}(y), T_{z_0}(y)) = \abs{y_0 - z_0}$, sehingga
pertanyaan 3 memberi
\[
d_\infty(y[y_0], y[z_0]) \leq \frac{\abs{y_0 - z_0}}{1 - Lh} .
\]

\textbf{16.} Pada tiap potongannya, pertanyaan 15 yang diterapkan dengan
nilai di ujung kirinya sebagai data awal membatasi simpangan di
ujung kanannya:
\[
d_\infty \leq \frac{1}{1 - 1/2}\,(\text{simpangan kiri})
= 2\,(\text{simpangan kiri}) .
\] Dengan
induksi atas $m$ potongannya, simpangan akhirnya paling banyak
$2^m\abs{y_0 - z_0}$, dan simpangan seragam pada seluruh
ruasnya menuruti batas yang sama (sebab supremum tiap potongan terkendali
pada tahapnya). Dengan potongan berpanjang $h \asymp \frac1{2L}$,
faktornya menjadi $2^m = 2^{\,\text{panjang}\cdot 2L}$: eksponensial
terhadap panjang intervalnya, yakni persis seperti diramalkan batas Gronwall
$\eu^{L\abs{t-t_0}}$, dengan konstanta yang lebih baik.

\textbf{17.} Skemanya sama: untuk $y \in X_h$,
\[
\abs{T_f(y)(t) - T_g(y)(t)} \leq \int_{t_0}^t \abs{f(s,y(s)) -
g(s,y(s))}\,\abs{\dd s} \leq \varepsilon h ,
\]
jadi pertanyaan 3 (dengan $T_f$ sebagai kontraksinya dan $T_g$ sebagai
pemetaan terusiknya) menghasilkan $d_\infty \leq \frac{\varepsilon h}{1 - Lh}$.

\textbf{18.} Menurut pertanyaan 17 yang diterapkan pada $f = f_\lambda$,
$g = f_\mu$: $d_\infty(y_\lambda, y_\mu) \leq \frac{Ch}{1 -
Lh}\abs{\lambda - \mu}$, jadi pemetaan penyelesaiannya bersifat \omterm{def:b2:metric:continuity}{Lipschitz}
dengan konstanta $\frac{Ch}{1-Lh}$.

\textbf{19.} Setiap hujahnya hanya memakai aksioma metrik,
kelengkapan panggungnya, batas $\abs{\int} \leq
\int\abs{\cdot}$, dan sifat \omterm{def:b2:metric:continuity}{Lipschitz} $f$ --- yang semuanya
tersedia bagi fungsi bernilai $\R^n$ dengan $d_\infty$ yang \omterm{def:b2:structures:generated}{dibangun} atas
salah satu jarak setara pada \cref{ex:b2:metric:examples}
(yang \omterm{def:b2:metric:complete}{lengkap} menurut \cref{thm:b2:metric:rncomplete}). Untuk $y' = Ay$
dengan $A$ yang nilpoten: $y^{(0)} \equiv (c_1, c_2)$,
\[
y^{(1)}(t) = (c_1, c_2) + \int_0^t (c_2, 0)\,\dd s
= (c_1 + tc_2,\; c_2),
\]
dan $Ay^{(1)}(s) = (c_2, 0)$ lagi: jadi $y^{(2)} = y^{(1)}$.
Iterasinya mandek mulai $n = 1$, tepat pada penyelesaian eksaknya
$y(t) = (c_1 + tc_2, c_2) = \eu^{tA}y(0)$ --- kenilpotenan
memangkas deret eksponensialnya, dan Picard menyadarinya.

\textbf{20.} Tetapkan $y \in X$ lalu misalkan $g_y(x) = y - \eta(x)$:
sebuah kontraksi-$k$ pada $X$ yang \omterm{def:b2:metric:complete}{lengkap}, jadi ada tepat satu $x$
dengan $x + \eta(x) = y$: karenanya pemetaan $\Phi = \mathrm{id} + \eta$
bersifat bijektif. Invers yang \omterm{def:b2:metric:continuity}{Lipschitz}: jika $\Phi(x) = y$ dan
$\Phi(x') = y'$, maka
\[
d(x, x') \leq d(y, y') + d\bigl(\eta(x), \eta(x')\bigr)
\leq d(y, y') + k\,d(x, x'),
\]
jadi $d(x, x') \leq \frac{1}{1-k}d(y, y')$. (Jaraknya di sini berasal
dari struktur normanya, sehingga $d(a - c, b - c) = d(a, b)$, dan itulah
yang dipakai ketaksamaan pertamanya.)

\textbf{21.} Pemetaan $g(x) = m + e\sin x$ bersifat Lipschitz-$e$ pada
$\R$ yang \omterm{def:b2:metric:complete}{lengkap} (lewat ketaksamaan nilai rata-rata, $\abs{g'} =
\abs{e\cos x} \leq e < 1$): jadi Banach memberi tepat satu penyelesaian
dan kekonvergenan global bagi iterasinya. Untuk $e = \frac12$, $m = 1$,
$x_0 = 1$: $x_1 = 1 + \frac12\sin 1 \approx 1.42074$, $d(x_1, x_0)
\approx 0.4207$, dan batas apriori
$\frac{(1/2)^n}{1/2}\cdot 0.4207 \leq 10^{-3}$ baru berlaku pada
$n = 10$: jadi sepuluh iterasi tersahkan (padahal nilai sejatinya $x
\approx 1.4987$ sebenarnya sudah tercapai sampai $10^{-3}$ sekitar $n =
5$).

\textbf{22.} Pakailah pemerian angkanya
(\cref{exo:b2:metric:8}): $C$ adalah himpunan semua jumlah
$\sum_{n\geq1}a_n3^{-n}$ dengan $a_n \in \{0, 2\}$. Maka $S_1(C) =
\{x/3 : x \in C\}$ adalah himpunan bagian yang $a_1 = 0$, dan $S_2(C) =
\{x/3 + 2/3\}$ himpunan bagian yang $a_1 = 2$: jadi gabungannya, atas
pilihan bebas $a_1$, tepat sama dengan $C$. Ketunggalannya dalam satu
kalimat: pada ruang himpunan bagian \omterm{def:b2:metric:compact}{kompak} tak kosong
$\intcc01$ yang dimetrikkan oleh jarak Hausdorff, pemetaan $A \mapsto
S_1(A) \cup S_2(A)$ merupakan kontraksi-$\frac13$ pada ruang yang
\omterm{def:b2:metric:complete}{lengkap}, jadi Banach hanya mengizinkan satu himpunan tetap --- dan jilid
Tahun ke-3 membuat metrik Hausdorff beserta hujah ini menjadi cermat.

\textbf{23.} Misalkan $Z = \{t \in J : y(t) = z(t)\}$: ia tak kosong
(sebab $t_0 \in Z$) dan tertutup di $J$ (sebab ia penyama dua pemetaan
\omterm{def:b2:metric:continuity}{kontinu}: prapeta $\{0\}$ di bawah $y - z$). \omterm{def:b2:metric:topology}{Terbuka}: jika $t_1 \in Z$,
terapkan teorema lokalnya (pertanyaan 8) di titik $(t_1, y(t_1))$,
pada persegi panjang tempat $f$ bersifat \omterm{def:b2:metric:continuity}{Lipschitz}: pada interval kecil
di sekitar $t_1$ baik $y$ maupun $z$ menyelesaikan masalah Cauchy yang
sama, jadi keduanya berimpit di sana (lewat ketunggalan tanpa syarat pada
pertanyaan 9): sehingga sebuah persekitaran $t_1$ berada di $Z$. Adapun
himpunan bagian tak kosong interval $J$ yang sekaligus \omterm{def:b2:metric:topology}{terbuka} dan
tertutup di $J$ pastilah seluruh $J$ (sebab interval bersifat \omterm{def:b2:metric:connected}{terhubung},
\cref{thm:b2:metric:connectedness}): jadi $y = z$ pada $J$.

\textbf{24.} Di sini $f(t, y) = \alpha(t)y + \beta(t)$ bersifat
Lipschitz-$L$ pada $y$ di seluruh $\intcc AB \times \R$ dengan $L =
\sup\abs\alpha$ (yang hingga, sebab $\alpha$ \omterm{def:b2:metric:continuity}{kontinu} pada sebuah ruas),
dan tak diperlukan pita $\intcc{y_0 - b}{y_0 + b}$ apa pun: ambil $X =
C(\intcc AB)$ seluruhnya, dan $T$ terdefinisi dengan baik di sana.
Induksi pertanyaan 8 berjalan kata demi kata dan memberi
$d_\infty(T^ny, T^nz) \leq \frac{(L(B - A))^n}{n!}d_\infty(y, z)$: jadi
suatu iteratnya merupakan kontraksi berapa pun panjang $B - A$, dan
pertanyaan 4 merampungkannya: ada satu dan hanya satu penyelesaian pada
seluruh ruas itu. Kelinearannya masuk tepat sekali: ia membuat batas
\omterm{def:b2:metric:continuity}{Lipschitznya} global pada $y$, sehingga persegi panjang beserta $h \leq
b/M$-nya tersingkir.

\textbf{25.} \emph{Kelengkapan} mengubah barisan Cauchy berisi
iterat menjadi penyelesaian yang sungguh ada (pertanyaan 1, 6, 8), dan
ketiadaannya membiarkan $\sqrt2$ lolos dari $\Q$ (pertanyaan 13).
\emph{Kontraksi} memberi ketunggalan, algoritme, dan
batang galatnya (pertanyaan 7, 14). \emph{Siasat iterat} menyingkirkan
syarat kekecilan $Lh < 1$, sehingga interval yang tersahkan
hanya bergantung pada $M$, bukan pada $L$ --- dan membuat persamaan
linear menjadi global (pertanyaan 8, 24). \emph{Keterhubungan} menaikkan
ketunggalan lokal menjadi ketunggalan global (pertanyaan 23). Adapun
contoh penyangkalnya: $2\sqrt{\abs y}$ menjaga syarat \omterm{def:b2:metric:continuity}{Lipschitz}
(pertanyaan 11), $y^2$ menjaga kelokalannya (pertanyaan 12), dan $\Q$
menjaga kelengkapannya (pertanyaan 13). Puncaknya adalah teorema
Picard--Lindelöf (pertanyaan 8); ketika $f$ sekadar
\omterm{def:b2:metric:continuity}{kontinu}, keberadaannya bertahan tetapi ketunggalannya mati, dan
buktinya menukar kontraksi dengan kekompakan himpunan fungsi
--- yakni teorema Peano lewat Arzelà--Ascoli, pada jilid Tahun ke-3.

\end{solution}
